%% starter: A problem set: numbered problems, each with a solution, and a figures folder
\documentclass[11pt]{article}
\usepackage[margin=2.5cm]{geometry}
\usepackage{amsmath}
\usepackage{graphicx}
\usepackage{siunitx}

% A problem is a numbered heading; its solution follows it.
\newcounter{problem}
\newcommand{\problem}{\stepcounter{problem}\section*{Problem \theproblem}}
\newenvironment{solution}{\par\medskip\noindent\textbf{Solution.}\ }{\par}

\begin{document}

\noindent\textbf{{{Course}}}\hfill\textbf{{{Title}}}\
{{Author}}\hfill\today
\par\medskip\hrule

\problem
% State the problem in your own words, then solve it.
A ball is thrown straight up at \qty{12}{m/s}. How high does it rise?

\begin{solution}
At the top the ball's speed is zero, so $v^2 = v_0^2 - 2gh$ gives
\begin{equation}
  h = \frac{v_0^2}{2g} = \frac{(\qty{12}{m/s})^2}{2(\qty{9.81}{m/s^2})} \approx \qty{7.3}{m}.
\end{equation}
\end{solution}

\problem
A second problem.

\begin{solution}
% align lines up the = signs; end each line but the last with \.
\begin{align}
  F &= ma \
    &= \qty{2.0}{kg} \times \qty{3.0}{m/s^2} = \qty{6.0}{N}.
\end{align}
\end{solution}

\problem
A third problem, with a sketch.

\begin{solution}
\begin{center}
  % Put your sketch in figures/ and uncomment the next line, naming it.
  % \includegraphics[width=0.6\linewidth]{figures/sketch.pdf}
  \fbox{\parbox[c][4cm][c]{0.6\linewidth}{\centering Your figure here}}
\end{center}
\end{solution}

\end{document}
